\documentclass[11pt,a4paper]{article}
\usepackage[french]{babel}
\usepackage[utf8]{inputenc}
\usepackage[T1]{fontenc}
\usepackage{geometry}
\usepackage{booktabs}
\usepackage{array}
\usepackage{xcolor}
\usepackage{titlesec}
\usepackage[hidelinks]{hyperref}
\usepackage{enumitem}
\usepackage{fancyhdr}
\usepackage{graphicx}
\usepackage{amsmath}
\usepackage{float}
\usepackage{caption}
\usepackage{subcaption}

\geometry{margin=2.5cm}

\definecolor{headerblue}{RGB}{31,73,125}

\pagestyle{fancy}
\fancyhf{}
\rhead{\textcolor{headerblue}{\small Rapport — Clustering des séries temporelles (FlexiMax)}}
\lhead{\small \leftmark}
\rfoot{\thepage}

\titleformat{\section}{\large\bfseries\color{headerblue}}{\thesection}{0.6em}{}[\titlerule]
\titleformat{\subsection}{\normalsize\bfseries\color{headerblue!80}}{\thesubsection}{0.6em}{}
\titleformat{\subsubsection}{\normalsize\bfseries\itshape}{\thesubsubsection}{0.6em}{}

% Adapter ce chemin si les figures sont ailleurs
\graphicspath{{figures/}}

\newcommand{\code}[1]{\texttt{#1}}

\begin{document}

\begin{titlepage}
    \centering
    \vspace*{2cm}
    {\LARGE\bfseries\color{headerblue} Clustering des séries temporelles de consommation énergétique\par}
    \vspace{0.5cm}
    {\large Identification de jours types de consommation électrique résidentielle\par}
    \vspace{1.5cm}
    {\large Projet FlexiMax --- MINES Paris--PSL / ARMINES\par}
    \vspace{0.3cm}
    {\normalsize Notebooks : \code{timeseries\_clustering.ipynb} et \code{timeseries\_clustering\_multi.ipynb}\par}
    \vfill
    {\normalsize Yassmin Zouarhi \par}
    {\normalsize \today\par}
\end{titlepage}

\tableofcontents
\newpage

\section{Objectif et démarche générale}

L'objectif de ce travail est d'identifier, à partir des séries temporelles de consommation
électrique des bâtiments résidentiels (données ResStock, pas de temps de 15 minutes), des
\textbf{jours types de consommation} par une approche de \textbf{clustering non supervisé}
(K-Means). Chaque journée de consommation est traitée comme une observation à part entière : les
séries annuelles sont restructurées en profils journaliers, normalisés, puis regroupés selon leurs
similarités.

Deux notebooks documentent cette démarche, dans un ordre de complexité croissante :

\begin{itemize}
    \item \textbf{\code{timeseries\_clustering.ipynb}} : mise au point de la méthode sur un
    bâtiment de référence (\code{347201-0}), puis validation sur un second bâtiment
    (\code{10672-0}) pour vérifier que les résultats ne sont pas un artefact d'un cas particulier.
    \item \textbf{\code{timeseries\_clustering\_multi.ipynb}} : généralisation de l'approche à
    \textbf{100 bâtiments tout-électriques}, avec construction d'un espace de features
    multivarié (forme du profil, amplitude énergétique, météo), analyse des corrélations, et
    mise en place d'une fonction de classification d'un nouveau jour.
\end{itemize}

La méthodologie commune aux deux notebooks se décompose en six étapes :

\begin{enumerate}
    \item Chargement des séries temporelles et sélection de la variable de consommation
    électrique totale (\code{out.electricity.total.energy\_consumption..kwh}).
    \item Transformation des séries annuelles en matrices \emph{jours~$\times$~pas de temps}.
    \item Normalisation des profils journaliers, selon deux logiques distinctes : la
    \emph{forme} (indépendante du niveau d'énergie) et l'\emph{amplitude} (niveau global de
    consommation).
    \item Détermination du nombre optimal de clusters $k$ (méthode du coude + score de
    silhouette).
    \item Application de l'algorithme K-Means.
    \item Analyse et visualisation des clusters obtenus (profils moyens, composition
    calendaire, projection PCA, etc.).
\end{enumerate}

\section{Notebook 1 --- \code{timeseries\_clustering.ipynb}}

\subsection{Chargement et exploration du bâtiment de référence}

Le bâtiment \code{347201-0} sert de cas d'étude principal. Les données brutes sont chargées
depuis \code{data/raw/347201-0.parquet}. Une première exploration visuelle décompose la
consommation totale par usage (colonnes \code{out.electricity.*.energy\_consumption..kwh}) et
affiche :

\begin{itemize}
    \item un graphique en aires empilées de la consommation journalière par poste
    (chauffage, climatisation, éclairage, électroménager, etc.), la colonne agrégée \code{total}
    étant retirée de cette décomposition ;
    \item une superposition de toutes les courbes journalières horaires, colorées selon
    semaine/weekend, afin d'avoir un premier aperçu visuel de la variabilité des profils.
\end{itemize}

\subsection{Extraction des profils journaliers}

La série annuelle (pas de temps de 15 minutes, 96 points/jour) est réorganisée en une matrice
\emph{jours~$\times$~96 pas de temps} : chaque ligne correspond à une journée complète de
consommation. Seuls les jours complets sont conservés (troncature du reste de la division
par 96), et une date représentative est associée à chaque ligne. Les journées contenant des
valeurs manquantes (\code{NaN}) sont ensuite supprimées.

\subsection{Normalisation : forme et amplitude}

Deux stratégies de normalisation sont appliquées en parallèle, car elles répondent à deux
questions différentes :

\begin{itemize}
    \item \textbf{Option A --- amplitude} (\code{StandardScaler} global) : standardisation
    colonne par colonne (pas de temps) sur l'ensemble des jours. Le \emph{niveau} d'énergie est
    conservé, ce qui permet de distinguer les journées de forte et de faible consommation.
    \item \textbf{Option B --- forme} (z-score intra-jour) : chaque journée est centrée-réduite
    par sa propre moyenne et son propre écart-type
    ($\frac{x - \bar{x}_{\text{jour}}}{\sigma_{\text{jour}}}$). Cette normalisation isole la
    \emph{forme} du profil (répartition horaire) indépendamment du niveau absolu d'énergie
    consommée.
\end{itemize}

Les deux jeux de données ($X$ = forme, $Y$ = amplitude) sont conservés pour la suite de
l'analyse, afin de comparer les deux points de vue.

\subsection{Choix du nombre de clusters}

Pour chaque normalisation, un nombre de clusters $k \in [2, 14]$ est testé avec K-Means
(\code{n\_init=10}). Pour chaque $k$ sont calculés :

\begin{itemize}
    \item l'\textbf{inertie} (somme des distances au carré aux centroïdes), utilisée par la
    méthode du coude (\code{KneeLocator}, courbe convexe décroissante) pour proposer un
    $k$ optimal ;
    \item le \textbf{score de silhouette} (calculé sur un sous-échantillon de 5000 jours pour
    accélérer le calcul), qui mesure la qualité de séparation des clusters.
\end{itemize}

\begin{figure}[H]
    \centering
    \includegraphics[width=0.9\textwidth]{kmeans_elbow_silhouette.png}
    \caption{Méthode du coude (inertie) et score de silhouette en fonction de $k$, pour les
    deux normalisations (forme et amplitude).}
    \label{fig:elbow-silhouette}
\end{figure}

Deux valeurs de $k$ sont ainsi obtenues : $k_{\text{forme}}$ et $k_{\text{amplitude}}$, chacune
utilisée pour un clustering K-Means final indépendant (\code{n\_init=20}).

\subsection{Clustering final et visualisation des jours types}

Pour chaque clustering (forme et amplitude), les résultats sont enrichis de métadonnées
calendaires (jour de la semaine, indicateur weekend, mois) puis visualisés sous plusieurs
angles~:

\begin{itemize}
    \item \textbf{profils moyens par cluster} : courbe moyenne $\pm 1\sigma$ pour chaque
    cluster, avec un graphique interactif (Plotly) permettant de basculer d'un cluster à
    l'autre ;
    \item \textbf{composition calendaire} : répartition semaine/weekend et répartition
    mensuelle (heatmap) au sein de chaque cluster, pour interpréter le sens physique des jours
    types obtenus (ex. cluster \og{}été climatisation\fg{} vs cluster \og{}hiver chauffage\fg{}) ;
    \item \textbf{frise temporelle} : projection de l'appartenance aux clusters le long de
    l'année, pour visualiser la saisonnalité de l'affectation ;
    \item \textbf{projection PCA 2D} : réduction de dimension (jusqu'à 10 composantes,
    seuil de 90\,\% de variance expliquée retenu comme repère) puis projection des deux premières
    composantes, colorée par cluster, pour visualiser la séparation des groupes dans un plan.
\end{itemize}

\begin{figure}[H]
    \centering
    \begin{subfigure}[b]{0.48\textwidth}
        \centering
        \includegraphics[width=\textwidth]{cluster_daily_profiles.png}
        \caption{Profils journaliers moyens par cluster.}
    \end{subfigure}
    \hfill
    \begin{subfigure}[b]{0.48\textwidth}
        \centering
        \includegraphics[width=\textwidth]{pca_2d_clusters.png}
        \caption{Projection PCA 2D colorée par cluster.}
    \end{subfigure}
    \caption{Jours types identifiés pour le bâtiment de référence \code{347201-0}.}
    \label{fig:cluster-daily-pca}
\end{figure}

\begin{figure}[H]
    \centering
    \includegraphics[width=0.9\textwidth]{cluster_calendar_composition.png}
    \caption{Composition calendaire des clusters : répartition semaine/weekend et répartition
    mensuelle.}
    \label{fig:cluster-calendar}
\end{figure}

\begin{figure}[H]
    \centering
    \includegraphics[width=0.9\textwidth]{cluster_timeline.png}
    \caption{Frise temporelle des affectations de cluster sur l'année.}
    \label{fig:cluster-timeline}
\end{figure}

\subsection{Validation sur un second bâtiment (\code{10672-0})}

Afin de vérifier que la méthode ne surapprend pas les spécificités du bâtiment
\code{347201-0}, l'intégralité du pipeline (extraction des jours, suppression des \code{NaN},
normalisation forme/amplitude, choix de $k$, clustering, visualisations, PCA) est reproduite à
l'identique sur le bâtiment \code{10672-0} --- le même bâtiment de référence que celui utilisé
dans le notebook multivarié (sélectionné automatiquement parmi les 100 bâtiments
tout-électriques). Les jours types obtenus présentent une structure cohérente avec ceux du
premier bâtiment (mêmes grandes familles de profils : jour \og{}chauffage\fg{}, jour
\og{}climatisation\fg{}, jour \og{}intersaison\fg{}), ce qui conforte la robustesse de
l'approche avant son passage à l'échelle multi-bâtiments.

\subsection{Extension multivariée mono-bâtiment}

En fin de notebook, une première extension multivariée est testée sur le bâtiment
\code{347201-0} : en plus de la forme du profil, chaque jour est enrichi de son amplitude
énergétique totale et de variables météo (température extérieure, humidité relative,
température humide, taux d'humidité), agrégées par jour. Après réduction de la forme par PCA,
les variables d'amplitude et de météo sont concaténées, un $k$ optimal est recherché et un
clustering K-Means final est appliqué. Les clusters obtenus sont analysés via des boxplots
(température et amplitude par cluster) et une heatmap résumant la moyenne de chaque variable
par cluster. Cette extension préfigure directement l'approche multi-bâtiments détaillée dans
le second notebook.

\section{Notebook 2 --- \code{timeseries\_clustering\_multi.ipynb}}

Ce second notebook généralise l'approche à un échantillon de bâtiments, dans le but de dégager
des jours types \emph{robustes}, non spécifiques à un seul logement, et d'étudier les facteurs
(météo, calendrier) qui expliquent la variabilité de consommation.

\subsection{Sélection de 100 bâtiments tout-électriques}

Le filtre déjà validé dans \code{extraction\_timeseries\_oedi.ipynb} est réutilisé :
maison individuelle à un étage, chauffage gainé électrique, climatisation centrale, sans
véhicule électrique, sans piscine, sans photovoltaïque, dans les zones climatiques ASHRAE
3A/4A/5A. Un critère supplémentaire \textbf{strict tout-électrique} est ajouté : consommation
nulle de gaz naturel, fioul et propane. Cela garantit qu'aucun usage énergétique n'est
\og{}caché\fg{} sur un autre vecteur énergétique (chauffage secondaire au gaz, chauffe-eau au
fioul, etc.), condition nécessaire pour analyser proprement la flexibilité de la demande
\emph{électrique}.

Parmi les bâtiments candidats respectant ces critères, 100 bâtiments sont tirés aléatoirement
(\code{random\_state=42}) et la liste est persistée dans
\code{data/processed/buildings\_all\_electric\_100.parquet} pour garantir la reproductibilité
d'une exécution à l'autre.

\subsection{Téléchargement des séries temporelles}

Les séries temporelles des 100 bâtiments sélectionnés sont téléchargées depuis le data lake
OEDI (Amazon S3, \emph{ResStock AMY2018, release 1}) et stockées dans
\code{data/raw/timeseries\_100\_all\_electric/}. Le téléchargement est \textbf{idempotent} :
seuls les fichiers manquants sont récupérés, ce qui permet de relancer le notebook sans
retélécharger l'ensemble des données à chaque fois.

\subsection{Extraction des features jour \texorpdfstring{$\times$}{x} bâtiment}

Contrairement au premier notebook (pas de 15 minutes, 96 points/jour), l'extraction est ici
faite au \textbf{pas horaire} (24 points/jour), afin de réduire la dimensionnalité et le coût de
calcul sur 100 bâtiments. Pour chaque bâtiment, une fonction \code{build\_day\_features}
calcule, pour chaque jour~:

\begin{itemize}
    \item la \textbf{forme} du profil (z-score intra-jour sur les 24 valeurs horaires) ;
    \item l'\textbf{amplitude} (somme journalière de la consommation, en kWh) ;
    \item des statistiques météo journalières (moyenne, minimum, maximum) pour chacune des
    quatre variables : température extérieure sèche, humidité relative, température humide,
    taux d'humidité.
\end{itemize}

Les résultats sont concaténés sur l'ensemble des 100 bâtiments dans une table de métadonnées
(\code{meta\_all}) et une matrice de formes (\code{shape\_all}).

\subsection{Clustering exploratoire sur un bâtiment de référence}

Avant de passer au multivarié, le pipeline forme/amplitude du premier notebook est rejoué sur
un bâtiment de référence (le premier bâtiment de l'échantillon trié par identifiant), à titre
de vérification de cohérence au pas horaire. Une visualisation interactive complémentaire est
introduite : au lieu de n'afficher que la courbe moyenne par cluster, \textbf{toutes} les
journées individuelles sont superposées, ce qui permet de juger visuellement de l'homogénéité
réelle de chaque cluster (et pas seulement de sa moyenne).

\subsection{Construction du jeu de données multivarié multi-bâtiments}

C'est le cœur méthodologique de ce notebook. Trois familles de variables sont combinées pour
décrire chaque jour, mais \textbf{pas de la même façon} :

\begin{itemize}
    \item \textbf{Forme} (24 points/jour) : réduite par \textbf{PCA} (jusqu'à 10 composantes),
    afin de condenser l'information de courbe tout en limitant la dimensionnalité ;
    \item \textbf{Amplitude} : transformée en $\log(1+x)$ puis standardisée \emph{par
    bâtiment} (z-score intra-bâtiment sur l'amplitude log-transformée). Cette étape est
    essentielle : elle isole le \emph{jour type} (journée relativement plus ou moins
    consommatrice pour \emph{ce} bâtiment) de la \emph{taille} du bâtiment (un grand logement
    consommant structurellement plus qu'un petit) ;
    \item \textbf{Météo} : standardisée globalement (\code{StandardScaler}), mais
    \textbf{sans PCA}, afin de rester directement interprétable dans les analyses de
    corrélation qui suivent.
\end{itemize}

Le vecteur de features final est la concaténation horizontale de la forme réduite par PCA et
des variables d'amplitude/météo brutes standardisées :
$$X_{\text{multi}} = \big[\, \text{PCA}(\text{forme}) \;\big\|\; \text{amplitude}_z \;\big\|\;
\text{météo}_{\text{std}} \,\big]$$

\begin{figure}[H]
    \centering
    \includegraphics[width=0.9\textwidth]{pca_variance.png}
    \caption{Variance expliquée par composante et variance cumulée de la PCA appliquée à la
    forme des profils journaliers (multi-bâtiments).}
    \label{fig:pca-variance}
\end{figure}

\subsection{Choix de $k$ et clustering final}

La méthode du coude et le score de silhouette sont à nouveau appliqués sur
$X_{\text{multi}}$ pour choisir un nombre de clusters $k_{\text{multi}}$, puis un K-Means final
(\code{n\_init=20}) est ajusté. Le résultat est projeté sur un plan interprétable
\emph{sans} recourir à une deuxième PCA : l'axe des abscisses est la première composante
principale de la forme (PC1), l'axe des ordonnées est directement la variable météo moyenne la
plus représentative (température extérieure).

\subsection{Vérification : absence de dominance par un seul bâtiment}

Un point de vigilance spécifique au passage au multi-bâtiments est vérifié explicitement : il
faut s'assurer qu'un cluster ne soit pas simplement \og{}un bâtiment particulier\fg{} plutôt
qu'un véritable \emph{jour type} partagé par plusieurs logements. Une table croisée
(bâtiment~$\times$~cluster, normalisée par ligne) est visualisée sous forme de heatmap pour
contrôler que chaque cluster est bien composé d'une diversité de bâtiments.

\subsection{Analyse des clusters multivariés}

Les clusters obtenus sont caractérisés selon plusieurs axes~:

\begin{itemize}
    \item boxplots de la température extérieure moyenne et de l'amplitude relative
    (z-score par bâtiment) par cluster ;
    \item répartition semaine/weekend par cluster ;
    \item heatmap résumant la moyenne de chaque variable météo et de l'amplitude par cluster,
    pour une vue d'ensemble synthétique ;
    \item profils de forme moyens par cluster (courbe horaire normalisée).
\end{itemize}

\begin{figure}[H]
    \centering
    \includegraphics[width=0.9\textwidth]{cluster_multi_analysis.png}
    \caption{Analyse des clusters multivariés : météo, amplitude relative et répartition
    semaine/weekend.}
    \label{fig:cluster-multi-analysis}
\end{figure}

\begin{figure}[H]
    \centering
    \begin{subfigure}[b]{0.48\textwidth}
        \centering
        \includegraphics[width=\textwidth]{cluster_multi_summary_heatmap.png}
        \caption{Heatmap résumé (moyenne des variables par cluster).}
    \end{subfigure}
    \hfill
    \begin{subfigure}[b]{0.48\textwidth}
        \centering
        \includegraphics[width=\textwidth]{cluster_multi_shape_profiles.png}
        \caption{Profils de forme moyens par cluster.}
    \end{subfigure}
    \caption{Synthèse du clustering multivarié multi-bâtiments.}
    \label{fig:cluster-multi-summary}
\end{figure}

Une visualisation interactive complémentaire superpose toutes les journées réelles (en kWh,
sans normalisation de forme) par cluster sur l'ensemble des 100 bâtiments, ainsi qu'une
inspection ciblée d'un bâtiment particulier (ses propres journées, colorées par cluster), pour
vérifier qu'un même logement peut effectivement basculer d'un jour type à l'autre selon les
conditions du jour.

\subsection{Analyse de corrélation}

Pour comprendre \emph{ce qui explique} la variation de consommation (indépendamment de la
taille du bâtiment, grâce au z-score par bâtiment), une matrice de corrélation est calculée
entre~:

\begin{itemize}
    \item les variables météo (température, humidité) ;
    \item l'amplitude relative (\code{amplitude\_z}) ;
    \item les variables calendaires (indicateur weekend, mois) ;
    \item les deux premières composantes principales de la forme (PC1, PC2).
\end{itemize}

Les corrélations avec l'amplitude sont ensuite triées par force d'association (valeur absolue
décroissante) et représentées en barres horizontales, afin d'identifier la ou les variables les
plus déterminantes. La variable météo la plus corrélée à l'amplitude est enfin représentée en
nuage de points interactif (Plotly) contre l'amplitude, colorée par cluster, avec une droite de
tendance (régression OLS).

\subsection{Classification d'un nouveau jour (fonction \og{}reverse\fg{})}

Une fonction \code{classify\_new\_day} est développée pour projeter un \emph{nouveau} profil
journalier (24 valeurs horaires de consommation, variables météo du jour, identifiant du
bâtiment) dans l'espace du clustering déjà ajusté, et prédire son cluster d'appartenance. Le
même enchaînement de transformations que lors de l'entraînement est reproduit à l'identique :

\begin{enumerate}
    \item calcul de la forme (z-score intra-jour) puis projection dans l'espace PCA
    \emph{déjà ajusté} (\code{pca\_shape.transform}, pas de nouveau \code{fit}) ;
    \item calcul de l'amplitude log-transformée, standardisée avec les statistiques
    (moyenne/écart-type) du bâtiment concerné si elles sont connues, sinon avec la moyenne des
    statistiques de tous les bâtiments ;
    \item standardisation des variables météo avec le \code{StandardScaler} déjà ajusté ;
    \item assemblage du vecteur final dans le même ordre que lors de l'entraînement, puis
    prédiction du cluster via \code{kmeans.predict}.
\end{enumerate}

Une fonction \code{classify\_new\_days\_batch} généralise ce traitement à un ensemble de jours.
Un test de cohérence est également effectué : un jour déjà connu du jeu d'entraînement est
reclassé par cette fonction, et le cluster prédit est comparé au cluster réel pour valider que
le pipeline de projection est fidèle à celui utilisé lors du clustering initial.

\subsection{Restitution accessible (daltonisme)}

Le graphique final des profils de forme moyens par cluster est décliné dans une version
accessible aux personnes daltoniennes~: palette catégorielle validée (distances perceptuelles
$\Delta E$ vérifiées pour les daltonismes protan/deutan/tritan, sans paires rouge/vert
adjacentes), combinée à un style de trait différent par cluster, de sorte que l'identité
d'un cluster ne repose jamais sur la couleur seule. Des étiquettes directes (\og{}C0\fg{},
\og{}C1\fg{}, \ldots) sont placées à l'extrémité de chaque courbe, en encre neutre.

\begin{figure}[H]
    \centering
    \includegraphics[width=0.75\textwidth]{cluster_shape_profiles_colorblind_safe.png}
    \caption{Profils de forme moyens par cluster --- version adaptée au daltonisme (palette et
    styles de trait redondants).}
    \label{fig:colorblind-safe}
\end{figure}

\section{Synthèse}

\begin{itemize}
    \item Le \textbf{notebook 1} établit et valide la méthode de clustering de jours types sur
    des bâtiments individuels (deux bâtiments testés), en distinguant systématiquement
    \emph{forme} (comportement horaire) et \emph{amplitude} (niveau de consommation), et en
    esquissant une première extension multivariée intégrant la météo.
    \item Le \textbf{notebook 2} généralise cette approche à 100 bâtiments tout-électriques, au
    pas horaire, avec un espace de features combinant forme (réduite par PCA), amplitude
    relative (log + z-score par bâtiment) et météo (standardisée, non réduite). Il ajoute des
    garde-fous méthodologiques propres au passage à l'échelle (vérification de non-dominance
    par bâtiment), une analyse des déterminants de la consommation (corrélations), une brique
    d'inférence pour classer de nouveaux jours, et une restitution graphique accessible.
\end{itemize}

Ces travaux fournissent une base pour la suite du stage : les jours types ainsi identifiés
(et leurs déterminants météo/calendaires) pourront être mobilisés pour l'analyse de la
flexibilité de la demande résidentielle (upgrades \code{dr\_001}--\code{dr\_005}).

\end{document}
